% Options for packages loaded elsewhere
% Options for packages loaded elsewhere
\PassOptionsToPackage{unicode}{hyperref}
\PassOptionsToPackage{hyphens}{url}
\PassOptionsToPackage{dvipsnames,svgnames,x11names}{xcolor}
%
\documentclass[
  letterpaper,
  DIV=11,
  numbers=noendperiod]{scrartcl}
\usepackage{xcolor}
\usepackage{amsmath,amssymb}
\setcounter{secnumdepth}{-\maxdimen} % remove section numbering
\usepackage{iftex}
\ifPDFTeX
  \usepackage[T1]{fontenc}
  \usepackage[utf8]{inputenc}
  \usepackage{textcomp} % provide euro and other symbols
\else % if luatex or xetex
  \usepackage{unicode-math} % this also loads fontspec
  \defaultfontfeatures{Scale=MatchLowercase}
  \defaultfontfeatures[\rmfamily]{Ligatures=TeX,Scale=1}
\fi
\usepackage{lmodern}
\ifPDFTeX\else
  % xetex/luatex font selection
\fi
% Use upquote if available, for straight quotes in verbatim environments
\IfFileExists{upquote.sty}{\usepackage{upquote}}{}
\IfFileExists{microtype.sty}{% use microtype if available
  \usepackage[]{microtype}
  \UseMicrotypeSet[protrusion]{basicmath} % disable protrusion for tt fonts
}{}
\makeatletter
\@ifundefined{KOMAClassName}{% if non-KOMA class
  \IfFileExists{parskip.sty}{%
    \usepackage{parskip}
  }{% else
    \setlength{\parindent}{0pt}
    \setlength{\parskip}{6pt plus 2pt minus 1pt}}
}{% if KOMA class
  \KOMAoptions{parskip=half}}
\makeatother
% Make \paragraph and \subparagraph free-standing
\makeatletter
\ifx\paragraph\undefined\else
  \let\oldparagraph\paragraph
  \renewcommand{\paragraph}{
    \@ifstar
      \xxxParagraphStar
      \xxxParagraphNoStar
  }
  \newcommand{\xxxParagraphStar}[1]{\oldparagraph*{#1}\mbox{}}
  \newcommand{\xxxParagraphNoStar}[1]{\oldparagraph{#1}\mbox{}}
\fi
\ifx\subparagraph\undefined\else
  \let\oldsubparagraph\subparagraph
  \renewcommand{\subparagraph}{
    \@ifstar
      \xxxSubParagraphStar
      \xxxSubParagraphNoStar
  }
  \newcommand{\xxxSubParagraphStar}[1]{\oldsubparagraph*{#1}\mbox{}}
  \newcommand{\xxxSubParagraphNoStar}[1]{\oldsubparagraph{#1}\mbox{}}
\fi
\makeatother

\usepackage{color}
\usepackage{fancyvrb}
\newcommand{\VerbBar}{|}
\newcommand{\VERB}{\Verb[commandchars=\\\{\}]}
\DefineVerbatimEnvironment{Highlighting}{Verbatim}{commandchars=\\\{\}}
% Add ',fontsize=\small' for more characters per line
\usepackage{framed}
\definecolor{shadecolor}{RGB}{241,243,245}
\newenvironment{Shaded}{\begin{snugshade}}{\end{snugshade}}
\newcommand{\AlertTok}[1]{\textcolor[rgb]{0.68,0.00,0.00}{#1}}
\newcommand{\AnnotationTok}[1]{\textcolor[rgb]{0.37,0.37,0.37}{#1}}
\newcommand{\AttributeTok}[1]{\textcolor[rgb]{0.40,0.45,0.13}{#1}}
\newcommand{\BaseNTok}[1]{\textcolor[rgb]{0.68,0.00,0.00}{#1}}
\newcommand{\BuiltInTok}[1]{\textcolor[rgb]{0.00,0.23,0.31}{#1}}
\newcommand{\CharTok}[1]{\textcolor[rgb]{0.13,0.47,0.30}{#1}}
\newcommand{\CommentTok}[1]{\textcolor[rgb]{0.37,0.37,0.37}{#1}}
\newcommand{\CommentVarTok}[1]{\textcolor[rgb]{0.37,0.37,0.37}{\textit{#1}}}
\newcommand{\ConstantTok}[1]{\textcolor[rgb]{0.56,0.35,0.01}{#1}}
\newcommand{\ControlFlowTok}[1]{\textcolor[rgb]{0.00,0.23,0.31}{\textbf{#1}}}
\newcommand{\DataTypeTok}[1]{\textcolor[rgb]{0.68,0.00,0.00}{#1}}
\newcommand{\DecValTok}[1]{\textcolor[rgb]{0.68,0.00,0.00}{#1}}
\newcommand{\DocumentationTok}[1]{\textcolor[rgb]{0.37,0.37,0.37}{\textit{#1}}}
\newcommand{\ErrorTok}[1]{\textcolor[rgb]{0.68,0.00,0.00}{#1}}
\newcommand{\ExtensionTok}[1]{\textcolor[rgb]{0.00,0.23,0.31}{#1}}
\newcommand{\FloatTok}[1]{\textcolor[rgb]{0.68,0.00,0.00}{#1}}
\newcommand{\FunctionTok}[1]{\textcolor[rgb]{0.28,0.35,0.67}{#1}}
\newcommand{\ImportTok}[1]{\textcolor[rgb]{0.00,0.46,0.62}{#1}}
\newcommand{\InformationTok}[1]{\textcolor[rgb]{0.37,0.37,0.37}{#1}}
\newcommand{\KeywordTok}[1]{\textcolor[rgb]{0.00,0.23,0.31}{\textbf{#1}}}
\newcommand{\NormalTok}[1]{\textcolor[rgb]{0.00,0.23,0.31}{#1}}
\newcommand{\OperatorTok}[1]{\textcolor[rgb]{0.37,0.37,0.37}{#1}}
\newcommand{\OtherTok}[1]{\textcolor[rgb]{0.00,0.23,0.31}{#1}}
\newcommand{\PreprocessorTok}[1]{\textcolor[rgb]{0.68,0.00,0.00}{#1}}
\newcommand{\RegionMarkerTok}[1]{\textcolor[rgb]{0.00,0.23,0.31}{#1}}
\newcommand{\SpecialCharTok}[1]{\textcolor[rgb]{0.37,0.37,0.37}{#1}}
\newcommand{\SpecialStringTok}[1]{\textcolor[rgb]{0.13,0.47,0.30}{#1}}
\newcommand{\StringTok}[1]{\textcolor[rgb]{0.13,0.47,0.30}{#1}}
\newcommand{\VariableTok}[1]{\textcolor[rgb]{0.07,0.07,0.07}{#1}}
\newcommand{\VerbatimStringTok}[1]{\textcolor[rgb]{0.13,0.47,0.30}{#1}}
\newcommand{\WarningTok}[1]{\textcolor[rgb]{0.37,0.37,0.37}{\textit{#1}}}

\usepackage{longtable,booktabs,array}
\newcounter{none} % for unnumbered tables
\usepackage{calc} % for calculating minipage widths
% Correct order of tables after \paragraph or \subparagraph
\usepackage{etoolbox}
\makeatletter
\patchcmd\longtable{\par}{\if@noskipsec\mbox{}\fi\par}{}{}
\makeatother
% Allow footnotes in longtable head/foot
\IfFileExists{footnotehyper.sty}{\usepackage{footnotehyper}}{\usepackage{footnote}}
\makesavenoteenv{longtable}
\usepackage{graphicx}
\makeatletter
\newsavebox\pandoc@box
\newcommand*\pandocbounded[1]{% scales image to fit in text height/width
  \sbox\pandoc@box{#1}%
  \Gscale@div\@tempa{\textheight}{\dimexpr\ht\pandoc@box+\dp\pandoc@box\relax}%
  \Gscale@div\@tempb{\linewidth}{\wd\pandoc@box}%
  \ifdim\@tempb\p@<\@tempa\p@\let\@tempa\@tempb\fi% select the smaller of both
  \ifdim\@tempa\p@<\p@\scalebox{\@tempa}{\usebox\pandoc@box}%
  \else\usebox{\pandoc@box}%
  \fi%
}
% Set default figure placement to htbp
\def\fps@figure{htbp}
\makeatother





\setlength{\emergencystretch}{3em} % prevent overfull lines

\providecommand{\tightlist}{%
  \setlength{\itemsep}{0pt}\setlength{\parskip}{0pt}}



 


\KOMAoption{captions}{tableheading}
\makeatletter
\@ifpackageloaded{caption}{}{\usepackage{caption}}
\AtBeginDocument{%
\ifdefined\contentsname
  \renewcommand*\contentsname{Table of contents}
\else
  \newcommand\contentsname{Table of contents}
\fi
\ifdefined\listfigurename
  \renewcommand*\listfigurename{List of Figures}
\else
  \newcommand\listfigurename{List of Figures}
\fi
\ifdefined\listtablename
  \renewcommand*\listtablename{List of Tables}
\else
  \newcommand\listtablename{List of Tables}
\fi
\ifdefined\figurename
  \renewcommand*\figurename{Figure}
\else
  \newcommand\figurename{Figure}
\fi
\ifdefined\tablename
  \renewcommand*\tablename{Table}
\else
  \newcommand\tablename{Table}
\fi
}
\@ifpackageloaded{float}{}{\usepackage{float}}
\floatstyle{ruled}
\@ifundefined{c@chapter}{\newfloat{codelisting}{h}{lop}}{\newfloat{codelisting}{h}{lop}[chapter]}
\floatname{codelisting}{Listing}
\newcommand*\listoflistings{\listof{codelisting}{List of Listings}}
\makeatother
\makeatletter
\makeatother
\makeatletter
\@ifpackageloaded{caption}{}{\usepackage{caption}}
\@ifpackageloaded{subcaption}{}{\usepackage{subcaption}}
\makeatother
\usepackage{bookmark}
\IfFileExists{xurl.sty}{\usepackage{xurl}}{} % add URL line breaks if available
\urlstyle{same}
% fallback for those not using the hyperref driver hyperxmp:
\makeatletter
\@ifundefined{xmpquote}{\newcommand{\xmpquote}[1]{#1}}{}
\makeatother
\hypersetup{
  pdftitle={LEVAK validation framework},
  pdfauthor={Jose F. Meneses-Echavez},
  colorlinks=true,
  linkcolor={blue},
  filecolor={Maroon},
  citecolor={Blue},
  urlcolor={Blue},
  pdfcreator={LaTeX via pandoc}}


\title{LEVAK validation framework}
\author{Jose F. Meneses-Echavez}
\date{2026-11-13}
\begin{document}
\maketitle


\subsection{Background}\label{background}

Insert a short background sketching the report

\subsection{Objective}\label{objective}

The main aim is to prospectively evaluate the performance of
\texttt{AIscreenR} in screening references at the title and abstract
level in the different categories and folders presented in the living
evidence map for vaccine research
\href{https://github.com/jmeneses90/LEVAK}{LEVAK}, using the screening
decision of content experts as the gold-standard comparative benchmark.

\subsection{Uploading the libraries and
packages}\label{uploading-the-libraries-and-packages}

Centralizing all packages, libraries and functions used for the
evaluation and reporting into a dedicated chunk.

\begin{Shaded}
\begin{Highlighting}[]
\CommentTok{\# Load packages and libraries}
\FunctionTok{library}\NormalTok{(rsconnect)}
\FunctionTok{library}\NormalTok{(AIscreenR)}
\FunctionTok{library}\NormalTok{(writexl)}
\FunctionTok{library}\NormalTok{(stringr)}
\FunctionTok{library}\NormalTok{(dplyr)}
\FunctionTok{library}\NormalTok{(synthesisr)}
\FunctionTok{library}\NormalTok{(tibble)}
\FunctionTok{library}\NormalTok{(revtools)}
\FunctionTok{library}\NormalTok{(flextable)}
\FunctionTok{library}\NormalTok{(tinytex)}
\FunctionTok{library}\NormalTok{(tidyr)}
\FunctionTok{library}\NormalTok{(ggplot2)}
\FunctionTok{library}\NormalTok{(gt)}
\FunctionTok{library}\NormalTok{(webshot2)}
\FunctionTok{library}\NormalTok{(packrat)}
\FunctionTok{library}\NormalTok{(rsconnect)}
\FunctionTok{library}\NormalTok{(chromote)}
\FunctionTok{library}\NormalTok{(readxl)}
\FunctionTok{library}\NormalTok{(future)}
\FunctionTok{library}\NormalTok{(readr)}

\CommentTok{\# install the most recent pakcage update}
\NormalTok{devtools}\SpecialCharTok{::}\FunctionTok{install\_github}\NormalTok{(}\StringTok{"MikkelVembye/AIscreenR"}\NormalTok{, }\AttributeTok{quiet =} \ConstantTok{TRUE}\NormalTok{)}
\end{Highlighting}
\end{Shaded}

\subsection{API key}\label{api-key}

\texttt{AIscreenR} uses an API key to connect to a Large Language Model
(LLM), in this case, the OpenAI API to connect to ChatGPT. After
obtaining your API key (which is private and confidential), you can save
it using the following code:

\begin{Shaded}
\begin{Highlighting}[]
\CommentTok{\# get\_api\_key() }
\end{Highlighting}
\end{Shaded}

\section{Influenza vaccines}\label{influenza-vaccines}

The validation methods are presented separately for the two main
categories in LEVAK: a) vaccine effectiveness and safety, and b) burden
of disease and risk groups.

\subsection{Vaccine effectiveness and
safety}\label{vaccine-effectiveness-and-safety}

\subsubsection{First evaluation}\label{first-evaluation}

The literature searches yielded a total of 5451 unique references:

\begin{itemize}
\tightlist
\item
  2023 (1733 references)
\item
  2024 (1553 references)
\item
  2025 (1731 references)
\item
  2026 (434 references)
\end{itemize}

References from 2026 are used in this valdiation.

Further details on the LEVAK methodological framework are available at:
{[}https://osf.io/pmut5/overview{]}.

To develop a robust prompt, we first evaluated \texttt{AIscreenR}'s
performance by screening a random sample of 200 references, from the 434
references identified in 2026.

During the second week of April 2026, two content experts (Birgitte
Klüwer and Anna Hayman Robertson) screened these references for
eligibility based on predefined selection criteria (available in the
LEVAK protocol).

\subsection{Importing the dataset}\label{importing-the-dataset}

Importing the dataset and preparing it for screening. The dataset is in
RIS format. We will read the RIS file into a dataframe and then rename
some columns for better clarity.

\begin{Shaded}
\begin{Highlighting}[]
\CommentTok{\# references\_inf\_valid\_1 \textless{}{-} read\_ris\_to\_dataframe("C:/Users/JOFM/OneDrive {-} Folkehelseinstituttet/Vaksinasjon av voksne og risikogrupper {-} LEVAK/Influensa\_LEVAK/Effekt\_\&\_sikkerhet\_Inf/Influensa\_AIScreenR\_validering/data/Referanser\_First200/Influensa\_AIScreenR\_1st200/Influensa\_effekt\_validering\_1st200.txt")}

\CommentTok{\# \# Rename U4 to "human\_code"}
\CommentTok{\# references\_inf\_valid\_1 \textless{}{-} references\_inf\_valid\_1 \%\textgreater{}\% \# nolint}
\CommentTok{\#   rename(human\_code = U4)}
\CommentTok{\#   Rename record\_id to "ID"}
\CommentTok{\# references\_inf\_valid\_1 \textless{}{-} references\_inf\_valid\_1 \%\textgreater{}\%}
\CommentTok{\#   rename(ID = record\_id)}
\CommentTok{\#   colnames(references\_inf\_valid\_1) \# checking the column names to ensure they are correctly renamed.}

\CommentTok{\# \# saving the references as a variable in the workspace}
\CommentTok{\# saveRDS(references\_inf\_valid\_1, "data/references\_inf\_valid\_1.rds")}
\end{Highlighting}
\end{Shaded}

\subsection{Prompt}\label{prompt}

Following a first draft of the prompt prepared by Jose Meneses and
approved by the content experts, we ran the screening process in
\texttt{AIscreenR} as follows:

\begin{Shaded}
\begin{Highlighting}[]
\CommentTok{\# prompt\_influensa\_valid\_1 \textless{}{-}"Role: You are an expert systematic review screener expert in vaccinology and epidemiology. Task: Evaluate whether a study meets the following eligibility criteria: 1. Inclusion Criteria: A study must meet all the following criteria to be included:}
\CommentTok{\# Population: Human participants of any age, health status, or vaccination history. Intervention: The study evaluates the effects, efficacy, or safety of any influenza vaccine. The study must report data specifically on influenza vaccines.}
\CommentTok{\# Comparison: No restrictions (any comparator or no comparator is acceptable). Eligible Study Designs: Experimental studies: Randomized controlled trials (RCTs), non{-}randomized studies, quasi{-}experimental.}
\CommentTok{\# Observational: Prospective/retrospective cohort, population{-}based, registry{-}based, case{-}control, and nested case{-}control. Secondary/Other: Systematic reviews, meta{-}analyses, and post{-}marketing safety evaluations. Include guidelines only if explicitly based on systematic reviews. }
\CommentTok{\# 2. Exclusion Criteria. Exclude the study if any of the following apply: Not conducted in humans; does not report results specific to influenza vaccines; cross{-}sectional studies; surveys or purely self{-}reported data without outcome validation;  qualitative studies; narrative (non{-}systematic) reviews; case reports, case series, editorials, commentaries, letters without original or synthesized data.}
\CommentTok{\# Evaluation Rule: Does the study meet all the inclusion criteria? Only include the study if every criterion is satisfied. Provide a rationale for your assessment"}

\CommentTok{\# writeLines(prompt\_influensa\_valid\_1, "prompts/prompt\_influensa\_valid\_1.txt")}
\end{Highlighting}
\end{Shaded}

\subsection{Running the screening
process}\label{running-the-screening-process}

\begin{Shaded}
\begin{Highlighting}[]
\CommentTok{\#influensa\_valid\_1\_screening \textless{}{-}  tabscreen\_gpt(}
  \CommentTok{\#data = references\_inf\_valid\_1,}
  \CommentTok{\#prompt = prompt\_influensa\_valid\_1,}
  \CommentTok{\#studyid = ID,}
  \CommentTok{\#title = title,}
  \CommentTok{\#abstract = abstract,}
  \CommentTok{\#model = "gpt{-}5{-}mini",}
  \CommentTok{\#reps = 1,}
  \CommentTok{\#decision\_description = TRUE}
\CommentTok{\#)}
\end{Highlighting}
\end{Shaded}

Saving the screening results

\begin{Shaded}
\begin{Highlighting}[]
\CommentTok{\# \#\# saving the screening file (RDS)}
\CommentTok{\# saveRDS(influensa\_valid\_1\_screening, "influensa\_valid\_1\_screening.rds")}

\CommentTok{\# \# \# loading the results file}
\CommentTok{\# influensa\_valid1\_screening \textless{}{-} readRDS("outputs/influensa\_valid\_1\_screening.rds")}

\CommentTok{\# \# \# convert the rds file to Excel for result analysis}
\CommentTok{\# write.xlsx(influensa\_valid1\_screening,"outputs/influensa\_valid1\_screening\_results.xlsx")}
\end{Highlighting}
\end{Shaded}

\subsection{Analysis}\label{analysis}

\subsubsection{Performance metrics: analyzing the screening
results}\label{performance-metrics-analyzing-the-screening-results}

\begin{Shaded}
\begin{Highlighting}[]
\CommentTok{\# \#  Analysis {-} screen\_analyzer }
\CommentTok{\# influensa\_valid1\_analysis \textless{}{-} screen\_analyzer(influensa\_valid1\_screening, human\_decision = human\_code, key\_result = FALSE ) }

\CommentTok{\# glimpse(influensa\_valid1\_analysis)}
\CommentTok{\# saveRDS(influensa\_valid1\_analysis, "outputs/influensa\_valid1\_analysis.rds") }
\CommentTok{\# influensa\_valid1\_analysis\_excel \textless{}{-} write.xlsx(influensa\_valid1\_analysis, "outputs/influensa\_valid1\_analysis\_excel.xlsx")}
\end{Highlighting}
\end{Shaded}

\subsubsection{Key results (validation
1)}\label{key-results-validation-1}

\begin{itemize}
\item
  \texttt{AIscreenR} included 81 studies and excluded 119 studies.
\item
  \texttt{AIscreenR} had zero false negatives, which means that it
  excluded no relevant study. The \textbf{recall (sensitivity)} is
  \textbf{1}, which means that it correctly identified all of the
  studies that were included by human coders.
\item
  \texttt{AIscreenR} included 43 studies (21\%) that human coders had
  excluded (\textbf{false positives}).
\item
  The \textbf{specificity} of \texttt{AIscreenR} is \textbf{73\%}, which
  means that it correctly excluded 73\% of the studies that were
  excluded by human coders. \texttt{AIscreenR} included 43 studies that
  human coders had excluded.
\item
  The NPV of 1 suggested that \texttt{AIscreenR} correctly excluded all
  the irrelevant studies (those excluded by human coders).
\item
  The overall agreement between human coders and \texttt{AIscreenR} was
  78\%, indicating a disagreement of 21\%.
\item
  The \textbf{precision}, the ability to include only references that
  should be included, was 47\%.
\end{itemize}

\paragraph{False positives analysis}\label{false-positives-analysis}

The LEVAK team is interested in understanding the reasons behind the
false positives. These are the studies that were included by
\texttt{AIscreenR} but excluded by human coders.

We will analyze the false positives to identify any common
characteristics or patterns that may have led to their inclusion by
\texttt{AIscreenR}. This analysis will help us understand the
limitations of \texttt{AIscreenR} and provide insights for improving its
performance in future screenings.

The LEVAK team used the insights from this false negative analysis to
refine the prompt and put it to the test in a second round of screening.

\subsubsection{Second evaluation}\label{second-evaluation}

The content experts (human\_coders) independently read the remaining 234
references from 2026. We used the same methodological approach as
described for the first evaluation above. The remaining 234 references
from 2026 are selected:

\begin{Shaded}
\begin{Highlighting}[]
\CommentTok{\# references\_inf\_valid\_2 \textless{}{-} read\_ris\_to\_dataframe("data/Influensa\_2d\_2026\_234\_validering.txt")}
\end{Highlighting}
\end{Shaded}

Renaming U4 to `human\_code':

\begin{Shaded}
\begin{Highlighting}[]
\CommentTok{\# references\_inf\_valid\_2 \textless{}{-} references\_inf\_valid\_2 \%\textgreater{}\% }
\CommentTok{\#   rename(human\_code = U4)}

\CommentTok{\#   \# Rename record\_id to "ID"}
\CommentTok{\# references\_inf\_valid\_2 \textless{}{-} references\_inf\_valid\_2 \%\textgreater{}\%}
\CommentTok{\#   rename(ID = record\_id)}
\CommentTok{\# colnames(references\_inf\_valid\_2) \# check the column names to ensure they are correctly renamed.}
\end{Highlighting}
\end{Shaded}

\subsection{Running the screening
process}\label{running-the-screening-process-1}

To preserve validity, we used the same prompt:

\begin{Shaded}
\begin{Highlighting}[]
\CommentTok{\# influensa\_valid\_2\_screening \textless{}{-}  tabscreen\_gpt(}
\CommentTok{\#   data = references\_inf\_valid\_2,}
\CommentTok{\#   prompt = prompt\_influensa\_valid\_1,}
\CommentTok{\#   studyid = ID,}
\CommentTok{\#   title = title,}
\CommentTok{\#   abstract = abstract,}
\CommentTok{\#   model = "gpt{-}5{-}mini",}
\CommentTok{\#   reps = 1,}
\CommentTok{\#   decision\_description = TRUE}
\CommentTok{\# )}

\CommentTok{\# saveRDS(influensa\_valid\_2\_screening, "influensa\_valid\_2\_screening.rds")}
\CommentTok{\# influensa\_valid\_2\_screening \textless{}{-} readRDS("influensa\_valid\_2\_screening.rds")}

\CommentTok{\# \# convert the rds file to Excel for result analysis (counting the number of included and excluded studies based on the AI screening results)}
\CommentTok{\# write.xlsx(influensa\_valid\_2\_screening,"influensa\_valid\_2\_screening\_results.xlsx")}
\end{Highlighting}
\end{Shaded}

\subsection{Analysis}\label{analysis-1}

\subsubsection{Performance metrics: analyzing the screening
results}\label{performance-metrics-analyzing-the-screening-results-1}

\begin{Shaded}
\begin{Highlighting}[]
\CommentTok{\#  Analysis {-} screen\_analyzer }
\CommentTok{\# influensa\_valid2\_analysis \textless{}{-} screen\_analyzer(influensa\_valid2\_screening, human\_decision = human\_code, key\_result = FALSE ) }
\CommentTok{\# glimpse(influensa\_valid2\_analysis)}
\CommentTok{\# saveRDS(influensa\_valid2\_analysis, "outputs/influensa\_valid2\_analysis.rds") }
\CommentTok{\# influensa\_valid2\_analysis\_excel \textless{}{-} write.xlsx(influensa\_valid2\_analysis, "outputs/influensa\_valid2\_analysis\_excel.xlsx")}
\end{Highlighting}
\end{Shaded}

\begin{Shaded}
\begin{Highlighting}[]
\DocumentationTok{\#\# creating the screening results file in Excel }
\CommentTok{\# influensa\_valid2\_screening\_results\_excel \textless{}{-} read\_excel("outputs/influensa\_valid2\_results\_excel.xlsx")}
\CommentTok{\# colnames(influensa\_valid2\_screening\_results\_excel)}


\CommentTok{\# \#\# puling out the false positives (FP) from the screening results }
\CommentTok{\# FP\_infvalid2 \textless{}{-} influensa\_valid2\_screening\_results\_excel\%\textgreater{}\% }
\CommentTok{\#   filter(answer\_data.human\_code == 0,}
\CommentTok{\#   answer\_data.decision\_binary == 1)}

\CommentTok{\# view(FP\_infvalid2) \# checking the file}
\CommentTok{\# write.xlsx(FP\_infvalid2,"outputs/FP\_infvalid2.xlsx") \# saving the file \{r\}}
\end{Highlighting}
\end{Shaded}

REVISAR LOS 4 human\_in\_gpt\_ex

REVISAR LOS ROW DATA PORQUE NO APARECEN EXCLUIDOS POR AMBOS

\textbf{ANALYSIS PREVIO}

\begin{Shaded}
\begin{Highlighting}[]
\CommentTok{\# \# Step 1. Count the number of included and excluded studies based on the AI screening results}
\CommentTok{\# influensa\_valid\_2\_screening$answer\_data\%\textgreater{}\%  }
\CommentTok{\#   count(decision\_binary)}

\CommentTok{\# \# Step 2. Analysis {-} screen\_analyzer}
\CommentTok{\# influensa\_valid2\_analysis \textless{}{-} screen\_analyzer(influensa\_valid\_2\_screening, human\_decision = human\_code) }
\CommentTok{\# view(influensa\_valid2\_analysis)}
\CommentTok{\# \# saving the analysis results}
\CommentTok{\# saveRDS(influensa\_valid2\_analysis, "influensa\_valid1\_analysis.rds")}
\end{Highlighting}
\end{Shaded}

\texttt{AIscreenR} included 77 studies and excluded 157 studies based on
the screening process.

\subsubsection{Confusion matrix}\label{confusion-matrix}

\subsubsection{Key results, second
validation}\label{key-results-second-validation}

\begin{itemize}
\item
  \texttt{AIscreenR} had four false negatives. Besides, its
  \textbf{recall (sensitivity) is 88\%}, which means that it correctly
  identified almost all of the studies that were included by human
  reviewers (nine out of ten references).
\item
  The \textbf{specificity of} \texttt{AIscreenR} is \textbf{76\%}, which
  means that it correctly excluded 76\% of the studies that were
  excluded by human reviewers.
\end{itemize}

\subsubsection{Third evaluation}\label{third-evaluation}

Immunogenicity is added into the prompt. We test the results comparing
against the previous two evaluations above.

\begin{Shaded}
\begin{Highlighting}[]
\CommentTok{\# prompt\_influensa\_valid\_2 \textless{}{-} "Role: You are an expert systematic review screener specializing in vaccinology and epidemiology. Task: Evaluate the study meets the eligibility criteria below, based strictly on its title and abstract.}
\CommentTok{\# Eligibility Criteria}
\CommentTok{\# 1. Inclusion Criteria (Must meet ALL of the following):}
\CommentTok{\# {-} Population: Human participants of any age, health status, or vaccination history.}
\CommentTok{\# {-} Intervention: Explicitly evaluates the effects, efficacy, safety or immunogenicity of influenza vaccines. The study MUST report specific outcome data for influenza vaccines. Include immunogenicity studies only if they explicitly address the antibody response.}
\CommentTok{\# {-} Comparator: Any comparator (e.g., no vaccine, placebo vaccine, another influenza vaccine, or vaccine against another disease).}
\CommentTok{\# {-} Eligible Study Designs:}
\CommentTok{\# b}
\end{Highlighting}
\end{Shaded}

\section{COVID-19 vaccines}\label{covid-19-vaccines}

\subsection{Vaccine effectiveness and
safety}\label{vaccine-effectiveness-and-safety-1}

Insert text

\subsection{Burden of disease and risk
groups}\label{burden-of-disease-and-risk-groups}

Insert text

\section{Pneumococcal vaccines}\label{pneumococcal-vaccines}

\subsection{Vaccine effectiveness and
safety}\label{vaccine-effectiveness-and-safety-2}

Importing the references

\begin{Shaded}
\begin{Highlighting}[]
\CommentTok{\# uploading the text file}
\CommentTok{\# ref\_PNK\_ES\_valid\_1 \textless{}{-} read\_ris\_to\_dataframe("validation/data/pnk/pnk\_es/PNK\_ef\_valid.txt")}
\CommentTok{\# colnames(ref\_PNK\_ES\_valid\_1)}
\end{Highlighting}
\end{Shaded}

Rename U4 to ``human\_code''

\begin{Shaded}
\begin{Highlighting}[]
\CommentTok{\# ref\_PNK\_ES\_valid\_1 \textless{}{-} ref\_PNK\_ES\_valid\_1 \%\textgreater{}\%}
\CommentTok{\#   rename(human\_code = U4)}

\CommentTok{\# \# checking that column "human\_code" exists}
\CommentTok{\# colnames(ref\_PNK\_ES\_valid\_1)}

\CommentTok{\# \# saving the references}
\CommentTok{\# write\_xlsx(ref\_PNK\_\_ES\_valid\_1, "validation/outputs/pnk/pnk\_es/ref\_PNK\_ES\_valid\_1.xlsx")}
\end{Highlighting}
\end{Shaded}

Data cleaning

\begin{Shaded}
\begin{Highlighting}[]
\CommentTok{\# \# keeping only variables needed for AIscreenR and relevant metadata}
\CommentTok{\# ref\_PNK\_ES\_valid\_1 \textless{}{-} ref\_PNK\_ES\_valid\_1 \%\textgreater{}\% }
\CommentTok{\#   select(abstract, keywords, author, doi, year, title, human\_code, eppi\_id, record\_id)}

\CommentTok{\# \# Count references missing titles}
\CommentTok{\# sum(is.na(ref\_PNK\_ES\_valid\_1$title) |}
\CommentTok{\#       str\_trim(ref\_PNK\_ES\_valid\_1$title) == "" |}
\CommentTok{\#       ref\_PNK\_ES\_valid\_1$title == "No information") \# 0 ref}

\CommentTok{\# \# Count references missing abstracts}
\CommentTok{\# sum(is.na(ref\_PNK\_ES\_valid\_1$abstract) |}
\CommentTok{\#       str\_trim(ref\_PNK\_ES\_valid\_1$abstract) == "" |}
\CommentTok{\#       ref\_PNK\_ES\_valid\_1$abstract == "No information") \# 8 ref}

\CommentTok{\# \# Count references missing BOTH title and abstract}
\CommentTok{\# sum(}
\CommentTok{\#   (is.na(ref\_PNK\_ES\_valid\_1$title) | str\_trim(ref\_PNK\_ES\_valid\_1$title) == "" | ref\_PNK\_ES\_valid\_1$title == "No information") \&}
\CommentTok{\#   (is.na(ref\_PNK\_ES\_valid\_1$abstract) | str\_trim(ref\_PNK\_ES\_valid\_1$abstract) == "" | ref\_PNK\_ES\_valid\_1$abstract == "No information")}
\CommentTok{\# ) \# 0 ref}
\end{Highlighting}
\end{Shaded}

Prompt \_ PNK ES

The prompt is optimized to specify that studies on children under 5 are
relevant only if they are at high risk due to underlying medical
conditions.

\begin{Shaded}
\begin{Highlighting}[]
\CommentTok{\# prompt\_PNK\_ES \textless{}{-} "Role: You are an expert systematic review screener in vaccinology and epidemiology.}
\CommentTok{\# Task: Evaluate whether a study meets the following eligibility criteria:}
\CommentTok{\# Inclusion Criteria: A study must meet all of the following criteria to be included:}
\CommentTok{\# • Population: Human participants aged 5 years or older, regardless of health status or vaccination history. Studies involving children younger than 5 years are included only if the children are at high risk of severe pneumococcal disease due to underlying medical conditions (comorbidities). }
\CommentTok{\# • Intervention: The study evaluates the effectiveness, efficacy or safety of any pneumococcal vaccine. The study must report data specifically on pneumococcal vaccines. }
\CommentTok{\# • Comparison: Any comparator (e.g., no vaccine, placebo vaccine, another pneumococcal vaccine, or vaccine against another disease).}
\CommentTok{\# • Eligible Study Designs: }
\CommentTok{\# Experimental studies: Randomized controlled trials (parallel, cross{-}over, cluster, factorial, adaptive, and N{-}of{-}1 trials) and quasi{-}experimental studies (non{-}randomized trials, interrupted{-}time series, and difference{-}in{-}differences).}
\CommentTok{\# Observational studies: Post{-}marketing surveillance (Phase IV trials), case{-}control studies (traditional, test{-}negative, indirect cohort/Broome method, nested case{-}control, self{-}controlled case series), case{-}crossover studies, and cohort studies (prospective and retrospective).}
\CommentTok{\# Health{-}economic studies: Cost{-}effectiveness analysis (CEA), cost{-}utility analysis (CUA), cost{-}minimization analysis (CMA), cost{-}benefit analysis (CBA), and cost{-}consequence analysis.}
\CommentTok{\# Secondary research: Systematic reviews (with or without meta{-}analysis), overviews of reviews, living evidence maps, and clinical guidelines explicitly based on systematic reviews.}
\CommentTok{\# 2.    Exclusion Criteria: Exclude the study if any of the following apply: }
\CommentTok{\# Missing Information: Insufficient details in the title/abstract to confirm eligibility; missing abstract (unless the title explicitly specifies an eligible study design); or missing title; }
\CommentTok{\# Out of scope: Studies reporting on routine childhood vaccination; studies on the indirect effects of childhood pneumococcal conjugate vaccination; studies not conducted in humans; studies that only report on immunogenicity; studies that only focus on either otitis media or nasopharyngeal carriage; ecological impact studies without individual{-}level data.}
\CommentTok{\# Ineligible study design: Cross{-}sectional studies; surveys or self{-}reported data lacking outcome validation; qualitative studies; narrative (non{-}systematic) reviews; scoping reviews; mapping reviews; case reports; case series; editorials; conference proceedings; commentaries; replies to editor; and letters without original data.}
\CommentTok{\# Constraints \& Rules}
\CommentTok{\# {-} Base your decision ONLY on the provided title, abstract, and keywords. Full{-}text evaluation is not possible to confirm eligibility.}
\CommentTok{\# {-} Include a study ONLY if every inclusion criterion is met and no exclusion criterion applies.}
\CommentTok{\# Required Output Format}
\CommentTok{\# {-} **Decision:** [Include / Exclude]}
\CommentTok{\# {-} **Primary Reason:** [Briefly state the core reason for the decision]}
\CommentTok{\# {-} **Rationale:** [Provide a concise, step{-}by{-}step breakdown referencing specific criteria]"}

\CommentTok{\# }\AlertTok{\#\#\#}\CommentTok{ saving the prompt}
\CommentTok{\# writeLines(prompt\_PNK\_ES, "validation/prompts/pnk/pnk\_es/prompt\_PNK\_ES.txt")}
\end{Highlighting}
\end{Shaded}

Screening command

\begin{Shaded}
\begin{Highlighting}[]
\CommentTok{\# PNK\_ES\_screening1 \textless{}{-} tabscreen\_gpt(}
\CommentTok{\#   data = ref\_PNK\_ES\_valid\_1,}
\CommentTok{\#   prompt = prompt\_PNK\_ES,}
\CommentTok{\#   studyid = record\_id,}
\CommentTok{\#   title = title,}
\CommentTok{\#   abstract = abstract,}
\CommentTok{\#   model = "gpt{-}5{-}mini",}
\CommentTok{\#   reps = 1,}
\CommentTok{\#   decision\_description = TRUE}
\CommentTok{\# )}

\CommentTok{\# \# saving the screening results}
\CommentTok{\# saveRDS(PNK\_ES\_screening1,"validation/outputs/pnk/pnk\_es/PNK\_ES\_screening1.rds")}

\CommentTok{\# \# saving the screening results as an Excel file}
\CommentTok{\# write\_xlsx(as.data.frame(PNK\_ES\_screening1), "validation/outputs/pnk/pnk\_es/PNK\_ES\_screening1.xlsx")}
\end{Highlighting}
\end{Shaded}

Analysing the screening results

Counting inclusion and exclusion from AIscreenR

\begin{Shaded}
\begin{Highlighting}[]
\CommentTok{\# \# Count the number of included and excluded studies from AIscreenR}
\CommentTok{\# PNK\_ES\_screening1$answer\_data\%\textgreater{}\% }
\CommentTok{\# count(decision\_binary)}
\end{Highlighting}
\end{Shaded}

\texttt{AIscreenR} included 61 studies, and excluded 189 studies.

Detailed analysis

\begin{Shaded}
\begin{Highlighting}[]
\CommentTok{\# screen\_analyzer(PNK\_ES\_screening1, human\_decision = human\_code, key\_result = TRUE)}
\CommentTok{\# print(screen\_analyzer(PNK\_ES\_screening1, human\_decision = human\_code, key\_result = FALSE), width = Inf)}


\CommentTok{\# PNK\_ES\_screening1\_results \textless{}{-} screen\_analyzer(}
\CommentTok{\#     PNK\_ES\_screening1,}
\CommentTok{\#     human\_decision = human\_code,}
\CommentTok{\#     key\_result = FALSE}
\CommentTok{\# )}
\CommentTok{\# print(PNK\_ES\_screening1\_results) \# visualising }

\CommentTok{\# saveRDS(PNK\_ES\_screening1\_results, "validation/outputs/pnk/pnk\_es/PNK\_ES\_screening1\_results.rds") \# saving the R object}
\end{Highlighting}
\end{Shaded}

Screen\_analyzer yields the following results: n references: 250
screened Model: gpt-5-mini Recall: 96\% Specificity: 85\% Precision:
44\% Negative predictive value: 99\% F2: 0.38 false negatives: 1 -
called as ``human\_in\_gpt\_ex'' false positives: 34 - called as
``human\_ex\_gpt\_in'' true positives: 27 - called as
``human\_in\_gpt\_in'' true negatives: 188 - called as
``human\_ex\_gpt\_ex''

\subsection{Burden of disease and risk
groups}\label{burden-of-disease-and-risk-groups-1}

Insert text

\section{Attitudes toward vaccination \& vaccine
coverage}\label{attitudes-toward-vaccination-vaccine-coverage}

As reported in our \href{https://osf.io/pmut5/overview}{LEVAK protocol},
this evidence pillar consist of two categories: a) Attitudes toward
vaccination b) Interventions to increase vaccine coverage/uptake

For details on the evidence mapped out within the categories, please
visit our \href{https://osf.io/pmut5/overview}{LEVAK protocol}.

Search stratgies are built up differenlty for primary studies than for
secondary research. The search terms specific to the interventions
prioritzed in LEVAK.

\subsubsection{Evidence synthesis on both attitudes toward vaccination
\& interventions to increase vaccine
coverage/uptake}\label{evidence-synthesis-on-both-attitudes-toward-vaccination-interventions-to-increase-vaccine-coverageuptake}

This section of the validation evaluates the performance of
\texttt{AIscreenR} in identifying the prioritized types of evidence
synthesis reporting on both estimates of public attitudes toward
vaccination and effect studies evaluating interventions to increase
vaccine coverage and impact public attitudes toward vaccination. In
other words, this validation captures secondary research on
cross-sectional studies describing estimates of public attitudes toward
vaccination and vaccine coverage, and then, in the effects section, it
covers the interventions that impact them.

\subsubsection{Attitudes toward vaccination \& vaccine coverage: primary
studies on effect of
interventions}\label{attitudes-toward-vaccination-vaccine-coverage-primary-studies-on-effect-of-interventions}

Date: 09.11.2026 - Ingvild run the searches only in MEDLINE without the
search terms of setting. This search yields 792 references. The
references were exported in RIS format and saved in the folder
``data/holdninger/hold\_tiltak\_prim'' with the name ``Test LEVAK tiltak
uten setting.txt''. Bo and Angela will check a random sample of between
50-100 references to see how many relevant references could have been
missed by this search.

\begin{Shaded}
\begin{Highlighting}[]
\CommentTok{\# library(AIscreenR)}
\CommentTok{\# testref \textless{}{-} read\_ris\_to\_dataframe("C:/Users/JOFM/OneDrive {-} Folkehelseinstituttet/Vaksinasjon av voksne og risikogrupper {-} LEVAK/LEVAK/validation/data/holdninger/hold\_tiltak\_prim/Test LEVAK tiltak uten setting.txt")}

\CommentTok{\# view(testref)}
\CommentTok{\# nrow(testref)}

\CommentTok{\# saveRDS(testref, "C:/Users/JOFM/OneDrive {-} Folkehelseinstituttet/Vaksinasjon av voksne og risikogrupper {-} LEVAK/LEVAK/validation/outputs/holdninger/hold\_tiltak\_prim/testref.rds")}
\CommentTok{\# write\_csv(testref, "outputs/holdninger/hold\_tiltak\_prim/testref.csv")}
\end{Highlighting}
\end{Shaded}





\end{document}
